\exercisesheader{}

% 13

\eoce{\qt{Peluang bersama dan peluang bersyarat\label{joint_cond}} P(A) = 0.3,
P(B) = 0.7
\begin{parts}
\item Dapatkah Anda menghitung P(A dan B) jika hanya mengetahui P(A) dan P(B)?
\item Dengan asumsi bahwa peristiwa A dan B muncul dari proses acak yang saling independen,
\begin{subparts}
\item berapakah P(A dan B)?
\item berapakah P(A atau B)?
\item berapakah P(A$|$B)?
\end{subparts}
\item Jika diketahui bahwa P(A dan B) = 0.1, apakah variabel acak yang menghasilkan
peristiwa A dan B saling independen?
\item Jika diketahui bahwa P(A dan B) = 0.1, berapakah P(A$|$B)?
\end{parts}
}{}

% 14

\eoce{\qt{Selai kacang \& jeli\label{pbj}} Misalkan 80\% orang menyukai selai
kacang, 89\% menyukai jeli, dan 78\% menyukai keduanya. Jika diketahui bahwa
seseorang yang dipilih secara acak menyukai selai kacang, berapakah peluang
orang tersebut juga menyukai jeli?
}{}

% 15

\eoce{\qt{Pemanasan global\label{global_warming}} Sebuah jajak pendapat Pew Research
menanyakan kepada 1.306 orang Amerika, ``Berdasarkan apa yang telah Anda baca dan
dengar, apakah ada bukti kuat bahwa suhu rata-rata bumi meningkat selama beberapa
dekade terakhir, atau tidak?'' Tabel berikut menunjukkan distribusi tanggapan menurut
partai dan ideologi; jumlah responden telah diganti dengan frekuensi relatif.
\footfullcite{globalWarming}
\begin{center}
\begin{tabular}{ll  ccc c}
                    &                           & \multicolumn{3}{c}{\textit{Tanggapan}} \\
\cline{3-5}
                    &                           & Bumi      & Tidak     & Tidak tahu    &   \\
                    &                           & memanas   & memanas   & Menolak       & Total\\
\cline{2-6}
                    & Republikan konservatif    & 0.11      & 0.20      & 0.02      & 0.33  \\
\textit{Partai dan} & Republikan moderat/liberal& 0.06      & 0.06      & 0.01      & 0.13 \\
\textit{ideologi}   & Demokrat moderat/konservatif& 0.25    & 0.07      & 0.02      & 0.34 \\
                    & Demokrat liberal          & 0.18      & 0.01      & 0.01      & 0.20\\
\cline{2-6}
                    &Total                      & 0.60      & 0.34      & 0.06      & 1.00
\end{tabular}
\end{center}
\begin{parts}
\item Apakah meyakini bahwa bumi sedang memanas dan menjadi seorang Demokrat liberal
merupakan peristiwa yang saling eksklusif?
\item Berapakah peluang responden yang dipilih secara acak meyakini bahwa bumi
sedang memanas atau merupakan seorang Demokrat liberal?
\item Berapakah peluang responden yang dipilih secara acak meyakini bahwa bumi
sedang memanas jika diketahui bahwa responden tersebut merupakan seorang Demokrat liberal?
\item Berapakah peluang responden yang dipilih secara acak meyakini bahwa bumi
sedang memanas jika diketahui bahwa responden tersebut merupakan seorang Republikan konservatif?
\item Apakah keyakinan responden tentang bumi sedang memanas tampak saling independen
dengan partai dan ideologinya? Jelaskan penalaran Anda.
\item Berapakah peluang responden yang dipilih secara acak merupakan seorang
Republikan moderat/liberal jika diketahui bahwa responden tersebut tidak meyakini
bahwa bumi sedang memanas?
\end{parts}
}{}

\D{\newpage}

% 16

\eoce{\qt{Jaminan kesehatan, frekuensi relatif\label{health_coverage_rel_freqs}}
Behavioral Risk Factor Surveillance System (BRFSS) adalah survei telepon tahunan
yang dirancang untuk mengidentifikasi faktor risiko pada populasi orang dewasa dan
melaporkan tren kesehatan yang sedang berkembang. Tabel berikut menampilkan distribusi
status kesehatan responden survei ini (prima, sangat baik, baik, cukup, buruk) dan
apakah mereka memiliki jaminan kesehatan.
\begin{center}
\begin{tabular}{rrrrrrrr}
& &  \multicolumn{5}{c}{\textit{Status kesehatan}} &  \\
\cline{3-7}
                    &       & Prima     & Sangat baik & Baik      & Cukup     & Buruk     & Total \\
\cline{2-8}
\textit{Jaminan}    & Tidak & 0.0230    & 0.0364      & 0.0427    & 0.0192    & 0.0050    & 0.1262 \\
\textit{kesehatan}  & Ya    & 0.2099    & 0.3123      & 0.2410    & 0.0817    & 0.0289    & 0.8738 \\
\cline{2-8}
                    & Total & 0.2329    & 0.3486      & 0.2838    & 0.1009    & 0.0338    & 1.0000
\end{tabular}
\end{center}
\begin{parts}
\item Apakah berstatus kesehatan prima dan memiliki jaminan kesehatan merupakan
peristiwa yang saling eksklusif?
\item Berapakah peluang individu yang dipilih secara acak berstatus kesehatan prima?
\item Berapakah peluang individu yang dipilih secara acak berstatus kesehatan prima
jika diketahui bahwa individu tersebut memiliki jaminan kesehatan?
\item Berapakah peluang individu yang dipilih secara acak berstatus kesehatan prima
jika diketahui bahwa individu tersebut tidak memiliki jaminan kesehatan?
\item Apakah berstatus kesehatan prima dan memiliki jaminan kesehatan tampak saling
independen?
\end{parts}
}{}

% 17

\eoce{\qt{Preferensi burger\label{burger_preferences}} Sebuah jajak pendapat SurveyUSA
tahun 2010 menanyakan kepada 500 penduduk Los Angeles, ``Manakah tempat hamburger
terbaik di California Selatan? Five Guys Burgers? In-N-Out Burger? Fat Burger?
Tommy's Hamburgers? Umami Burger? Atau tempat lain?'' Distribusi tanggapan menurut
jenis kelamin ditunjukkan berikut ini. \footfullcite{burgers}
\begin{center}
\begin{tabular}{l p{4cm} r r r }
                    &                       & \multicolumn{2}{c}{\textit{Jenis kelamin}} \\
\cline{3-4}
                    &                       & Laki-laki  & Perempuan    & Total \\
\cline{2-5}
                    & Five Guys Burgers     & 5     & 6         & 11 \\
                    & In-N-Out Burger       & 162   & 181       & 343 \\
\textit{Tempat}     & Fat Burger            & 10    & 12        & 22 \\
\textit{hamburger}  & Tommy's Hamburgers    & 27    & 27        & 54 \\
\textit{terbaik}    & Umami Burger          & 5     & 1         & 6 \\
                    & Lainnya               & 26    & 20        & 46 \\
                    & Tidak yakin           & 13    & 5         & 18 \\
\cline{2-5}
                    & Total                 & 248   & 252       & 500
\end{tabular}
\end{center}
\begin{parts}
\item Apakah menjadi perempuan dan paling menyukai Five Guys Burgers merupakan
peristiwa yang saling eksklusif?
\item Berapakah peluang laki-laki yang dipilih secara acak paling menyukai In-N-Out?
\item Berapakah peluang perempuan yang dipilih secara acak paling menyukai In-N-Out?
\item Berapakah peluang seorang laki-laki dan seorang perempuan yang berpacaran
keduanya paling menyukai In-N-Out? Nyatakan setiap asumsi yang Anda buat dan nilai
apakah menurut Anda asumsi tersebut wajar.
\item Berapakah peluang seseorang yang dipilih secara acak paling menyukai Umami
atau orang tersebut perempuan?
\end{parts}
}{}

\D{\newpage}

% 18

\eoce{\qt{Perkawinan asortatif\label{assortative_mating}} Perkawinan asortatif adalah
pola perkawinan nonacak ketika individu dengan genotipe dan/atau fenotipe serupa
berpasangan lebih sering daripada yang diharapkan dalam pola perkawinan acak.
Peneliti yang mengkaji topik ini mengumpulkan data warna mata 204 laki-laki
Skandinavia dan pasangan perempuan mereka. Tabel berikut merangkum hasilnya.
\footfullcite{Laeng:2007}
\begin{center}
\begin{tabular}{ll  ccc c}
                                        &           & \multicolumn{3}{c}{\textit{Pasangan (perempuan)}} \\
\cline{3-5}
                                        &           & Biru  & Cokelat   & Hijau     & Total \\
\cline{2-6}
                                        & Biru      & 78    & 23        & 13        & 114 \\
\multirow{2}{*}{\textit{Diri (laki-laki)}}& Cokelat & 19    & 23        & 12        & 54 \\
                                        & Hijau     & 11    & 9         & 16        & 36 \\
\cline{2-6}
                                        & Total     & 108   & 55        & 41        & 204
\end{tabular}
\end{center}
\begin{parts}
\item Berapakah peluang responden laki-laki yang dipilih secara acak atau
pasangannya memiliki mata biru?
\item Berapakah peluang responden laki-laki bermata biru yang dipilih secara acak
memiliki pasangan bermata biru?
\item Berapakah peluang responden laki-laki bermata cokelat yang dipilih secara
acak memiliki pasangan bermata biru? Bagaimana dengan peluang responden laki-laki
bermata hijau yang dipilih secara acak memiliki pasangan bermata biru?
\item Apakah warna mata responden laki-laki dan pasangannya tampak saling independen?
Jelaskan penalaran Anda.
\end{parts}
}{}

% 19

\eoce{\qt{Membuat diagram kotak\label{tree_drawing_box_plots}} Setelah mengikuti mata
kuliah pengantar statistika, 80\% mahasiswa dapat membuat diagram kotak dengan benar.
Di antara mahasiswa yang dapat membuat diagram kotak, 86\% lulus, sedangkan hanya
65\% mahasiswa yang tidak dapat membuat diagram kotak yang lulus.
\begin{parts}
\item Buat diagram pohon untuk skenario ini.
\item Hitung peluang seorang mahasiswa dapat membuat diagram kotak jika diketahui
bahwa mahasiswa tersebut lulus.
\end{parts}
}{}

% 20

\eoce{\qt{Predisposisi trombosis\label{tree_thrombosis}} Sebuah tes genetik digunakan
untuk menentukan apakah seseorang memiliki predisposisi terhadap \textit{trombosis},
yaitu pembentukan bekuan darah di dalam pembuluh darah yang menghambat aliran darah
melalui sistem peredaran darah. Diperkirakan bahwa 3\% orang benar-benar memiliki
predisposisi ini. Tes genetik tersebut akurat 99\% jika seseorang benar-benar memiliki
predisposisi, yang berarti peluang hasil tes positif ketika seseorang benar-benar
memiliki predisposisi adalah 0.99. Tes tersebut akurat 98\% jika seseorang tidak
memiliki predisposisi. Berapakah peluang seseorang yang dipilih secara acak dan
mendapat hasil tes positif untuk predisposisi tersebut benar-benar memiliki predisposisi?
}{}

% 21

\eoce{\qt{Penyebabnya tak pernah lupus\label{tree_lupus}} Lupus adalah fenomena medis ketika antibodi
yang seharusnya menyerang sel asing untuk mencegah infeksi justru menganggap protein
plasma sebagai benda asing, sehingga menimbulkan risiko tinggi terjadinya pembekuan
darah. Diperkirakan bahwa 2\% populasi menderita penyakit ini. Tes tersebut akurat
98\% jika seseorang benar-benar menderita penyakit ini. Tes tersebut akurat 74\%
jika seseorang tidak menderita penyakit ini. Ada sebuah kalimat dari acara televisi
Fox \emph{House} yang sering diucapkan setelah seorang pasien mendapat hasil tes
positif untuk lupus: ``Penyebabnya tak pernah lupus.'' Menurut Anda, adakah kebenaran dalam pernyataan
ini? Gunakan peluang yang sesuai untuk mendukung jawaban Anda.
}{}

% 22

\eoce{\qt{Jajak pendapat seusai pemungutan suara\label{tree_exit_poll}} Edison Research
mengumpulkan hasil jajak pendapat seusai pemungutan suara dari beberapa sumber untuk
pemilihan pencabutan mandat Scott Walker di Wisconsin. Mereka menemukan bahwa 53\% responden
memilih Scott Walker. Selain itu, mereka memperkirakan bahwa di antara responden yang
memilih Scott Walker, 37\% memiliki gelar dari perguruan tinggi, sedangkan 44\% responden
yang tidak memilih Scott Walker memiliki gelar dari perguruan tinggi. Misalkan kita memilih
secara acak seseorang yang mengikuti jajak pendapat tersebut dan mengetahui bahwa
orang itu memiliki gelar dari perguruan tinggi. Berapakah peluang bahwa orang tersebut
memilih Scott Walker?
\footfullcite{data:scott}
}{}
